\begin{proof}[Proof of \cref{obj:lem:averaged-threshold-process-bounded}]
\noindent\textbf{1. Continuity of the cell threshold.}
For \(u,u'\in\mathbb R_+^4\), write
\[
s_a(u)=u_{a0}+u_{a1}\quad(a\in\{0,1\}),
\qquad
s(u)=s_0(u)+s_1(u)=\sum_{a,y\in\{0,1\}}u_{ay}.
\]
Assign the quotient in \cref{obj:def:dual-process} the value \(g_a(0)=0\) at its zero denominator. The bounds below show that this is its continuous extension on the nonnegative cone. Since \(u\) is nonnegative, \(s(u)=0\) implies \(u=0\). When \(s(u)>0\), its two branches are \(s(u)u_{a1}/s_a(u)\) and \(u_{a1}/\epsilon\), according as \(s_a(u)\geq\epsilon s(u)\) or \(s_a(u)<\epsilon s(u)\). On each branch, its absolute coordinate derivatives are at most \(1+\epsilon^{-1}\); at zero, \(0\leq g_a(u)\leq s(u)\). Comparing along the segment between two nonnegative vectors therefore gives, for \(a\in\{0,1\}\),
\[
|g_a(u)-g_a(u')|
\leq(1+\epsilon^{-1})\sum_{\zeta\in\mathcal Z}|u_\zeta-u'_\zeta|.
\]
The positive-part map is \(1\)-Lipschitz. Thus, for \(\lambda\in[0,1]\),
\[
\begin{aligned}
|f_\lambda(u)-f_\lambda(u')|
&\leq 2|g_0(u)-g_0(u')|
   +|g_1(u)-g_1(u')|
   +\lambda|s(u)-s(u')|\\
&\leq \bigl[3(1+\epsilon^{-1})+1\bigr]
       \sum_{\zeta\in\mathcal Z}|u_\zeta-u'_\zeta|.
\end{aligned}
\]
Hence \(f_\lambda\) is continuous on the nonnegative cone for every fixed \(\lambda\in[0,1]\).
% lean: CausalSmith.Stat.DiscreteBudgetvalueCurve.thresholdFunReal_continuousOn_nonnegative

\noindent\textbf{2. The pilot rectangle.}
Fix any pilot counts and cell \(j\). The quantities in \cref{obj:def:jf-estimator} satisfy
\[
m=\frac n8>0,\qquad L=\log(ed)\geq1,\qquad
\delta=\frac Lm>0,\qquad h_{\zeta j}>0.
\]
For every \(\boldsymbol z\in[-1,1]^{\mathcal Z}\) and \(\zeta\in\mathcal Z\),
\[
0\leq\ell_{\zeta j}
\leq c^0_{\zeta j}+R_{\zeta j}z_\zeta
\leq u_{\zeta j}.
\]
The map \(\boldsymbol z\mapsto c_j^0+R_j\odot\boldsymbol z\) is affine and continuous. Since \(f_\lambda\) is continuous on the nonnegative cone, its composition with this map is continuous on \([-1,1]^{\mathcal Z}\) for every \(\lambda\in[0,1]\). This includes pilot counts arising when \(M=0\).
% lean: CausalSmith.Stat.DiscreteBudgetvalueCurve.pilotRectangle_threshold_pullback_continuousOn

\noindent\textbf{3. Continuity of each clipped cell estimate.}
Keep the counts fixed. The Jackson order and degree in \cref{obj:def:jf-estimator} satisfy \(K\geq2\) and \(D=2(K-1)\geq2\). For \(r=0,\ldots,D\), set
\[
\xi_r=\frac{r}{D+1},
\qquad
\boldsymbol\xi_\gamma=(\xi_{\gamma_\zeta})_{\zeta\in\mathcal Z},
\qquad
\vartheta_{\gamma,\zeta}=\arccos(\xi_{\gamma_\zeta}),
\quad
\gamma\in\{0,\ldots,D\}^{\mathcal Z}.
\]
At a fixed grid point, the defining Jackson integral and the threshold formula give
\[
\begin{aligned}
&\left|
P_{j,\lambda}(c_j^0+R_j\odot\boldsymbol\xi_\gamma)
-P_{j,\mu}(c_j^0+R_j\odot\boldsymbol\xi_\gamma)
\right|\\
&\qquad\leq |\lambda-\mu|
\int_{[-\pi,\pi]^{\mathcal Z}}
s\!\left(c_j^0+R_j\odot
\cos(\boldsymbol\vartheta_\gamma+\boldsymbol t)\right)
\prod_{\zeta\in\mathcal Z}J_K(t_\zeta)\,d\boldsymbol t,
\qquad \lambda,\mu\in[0,1],
\end{aligned}
\]
where cosine acts coordinatewise. The integral is finite: its argument lies in the pilot rectangle, and the integrand is continuous on a compact box. Also,
\[
|f_\lambda(c_j^0)-f_\mu(c_j^0)|
\leq |\lambda-\mu|\,s(c_j^0).
\]

For \(\alpha,\gamma\in\{0,\ldots,D\}^{\mathcal Z}\), define the fixed tensor interpolation weights by
\[
\boldsymbol z^\alpha:=\prod_{\zeta\in\mathcal Z}z_\zeta^{\alpha_\zeta},
\qquad
w_{\alpha,\gamma}:=
[\boldsymbol z^\alpha]\,
\prod_{\zeta\in\mathcal Z}
\prod_{\substack{0\leq r\leq D\\r\ne\gamma_\zeta}}
\frac{z_\zeta-\xi_r}{\xi_{\gamma_\zeta}-\xi_r},
\]
where \([\boldsymbol z^\alpha]\) extracts the indicated coefficient. Tensor Lagrange interpolation of the degree-\(D\) polynomial yields
\[
\beta_{j,\alpha}(\lambda)
=
\sum_{\gamma\in\{0,\ldots,D\}^{\mathcal Z}}
w_{\alpha,\gamma}
\left[
P_{j,\lambda}(c_j^0+R_j\odot\boldsymbol\xi_\gamma)
-f_\lambda(c_j^0)
\right].
\]
The preceding bounds show that every \(\beta_{j,\alpha}(\lambda)\) is continuous in \(\lambda\). All count-dependent factors in
\[
Z_j(\lambda)=f_\lambda(c_j^0)+
\sum_{\alpha\in\{0,\ldots,D\}^{\mathcal Z}}
\beta_{j,\alpha}(\lambda)
\prod_{\zeta\in\mathcal Z}
\frac{U_{\alpha_\zeta}(N_{\zeta j};c^0_{\zeta j})}
     {R_{\zeta j}^{\alpha_\zeta}}
\]
are fixed and have positive radii in their denominators. Thus \(Z_j\) is continuous. Finally,
\[
T_j(\lambda)=f_\lambda(c_j^0)+
\min\!\left\{d^{1/4}S_j,\,
\max\!\left\{-d^{1/4}S_j,\,
Z_j(\lambda)-f_\lambda(c_j^0)\right\}\right\}
\]
is continuous, since finite minima and maxima preserve continuity.
% lean: CausalSmith.Stat.DiscreteBudgetvalueCurve.jacksonCellStatistic_continuous_time

\noindent\textbf{4. Averaging and boundedness.}
Write \(\mathfrak S_n\) for the permutations of the \(n\) record positions. Conditional on the fixed sample, the average is the finite sum
\[
\widehat F(\lambda)
=
\sum_{M=0}^{n}
\frac{e^{-n/4}(n/4)^M}{M!}\,
\frac{1}{n!\,2^n}
\sum_{\sigma\in\mathfrak S_n}
\sum_{\omega\in\{0,1\}^n}
\left.\left(\sum_{j=1}^{d}T_j(\lambda)\right)
\right|_{M,\sigma,\omega}.
\]
Only the first \(M\) marks affect the count table; each marking of those retained records occurs \(2^{n-M}\) times in the displayed sum. Thus the \(2^n\) normalization gives the same average as \(2^M\) markings of the retained records. The vertical bar means that the counts are computed from that permutation and marking. Each cell summand \(T_j(\lambda)\) is continuous in \(\lambda\). The capped value for \(M>n\) is the constant zero. Hence \(\widehat F\) is continuous on \([0,1]\). Compactness of that interval gives the finite real number
\[
B=\max_{\lambda\in[0,1]}|\widehat F(\lambda)|,
\qquad
|\widehat F(\lambda)|\leq B
\quad\text{for every }\lambda\in[0,1].
\]
% lean: CausalSmith.Stat.DiscreteBudgetvalueCurve.averagedThresholdProcess_bounded
\end{proof}
